\documentclass[a4paper,11pt,twoside]{MECOM}

%%%%%%%%%%%%%%%%%%%%%%%%%%%%%%%%%%%%%%%%%%%%%%%%%%%%%%%%%%%%%%%%%%%%%%%%%%%
%               Parametros principales del documento                      %
%%%%%%%%%%%%%%%%%%%%%%%%%%%%%%%%%%%%%%%%%%%%%%%%%%%%%%%%%%%%%%%%%%%%%%%%%%%

% Titulo
\titulo{Implementaci\'on de la clase MECOM para generaci\'on de informes t\'ecnicos en \LaTeX{}}

% Autores
\autores{E. O. Fogliatto}{}

% Revisores
\revisores{}{}{}

% Revision de calidad
\calidad{A. Gonz\'alez}

% Aprobacion
\aprobacion{E. Dari}

%Objetivo
\objetivo{El presente documento tiene como objetivo mostrar el uso de la clase MECOM.cls para gererar informes t\'ecnicos.}

% Alcance
\alcance{Su utilizaci\'on est\'a restringida al Departamento de Mec\'anica Computacional. GOOD SAVE \LaTeX{}!}

% Numero de informe tecnico
\numeroIT{XX/20XX}

% Metadatos para pdf
\hypersetup{
    pdfauthor={Ezequiel O. Fogliatto},
    pdftitle={Implementaci\'on de la clase MECOM para generaci\'on de informes t\'ecnicos en \LaTeX{}},
    pdfkeywords={\LaTeX{}, Informes T\'ecnicos, Mecom},
    pdfcreator={CNEA},
    pdfsubject={Departamento de Mec\'anica Computacional, IN-ATN40MC-XX/20XX}    
    }
    
% Autores de revisiones
\definechangesauthor[name={Johann Sebastian Mastropiero}, color=blue]{mas}






\begin{document}

    % Creacion de la caratula
    \portada
    
    % Creacion del indice
    \tableofcontents   
    
    % Comienzo del desarrollo del documento
    \printnomenclature[2cm]


        
    \section{Uso de MECOM.cls}
    La clase {\bf MECOM.cls} est\'a basada en {\bf article} \cite{hefferon_minutes_2005}, y est\'a pensada para escribir informes t\'ecnicos de Mecom. Se recomienda consultar el archivo ITec\textunderscore Mecom.tex para ver la implementaci\'on de secciones como portada, \'indice, abreviaturas, referencias, etc.
    \nomenclature[A]{Mecom}{Mec\'anica Computacional}



    
    \subsection{Portada}
        El comando {\bf \textbackslash portada} genera la portada en una p\'agina nueva. Para completarla es necesario llenar los siguientes campos.
        \begin{itemize}
            \item titulo
            \item autores
            \item revisores
            \item calidad
            \item aprobacion
            \item objetivo
            \item alcance
            \item numeroIT (s\'olo la parte final, es decir n\'umero/a\~no)
        \end{itemize}
        \par
        Es necesario conservar el logo de CNEA\nomenclature[A]{CNEA}{Comisi\'on Nacional de Energ\'ia At\'omica} para poder compilar el documento. Junto con este documento se adjunta un archivo \emph{Makefile} para facilitar la compilaci\'on completa (\'indices, abreviaturas, y referencias).
        \par
        Los tama\~nos de letra de la portada se encuentran fijos en Arial 11pt. Al modificar el tama\~no global, la portada no se modifica.


        
    \subsection{\'Indice}
        Se redefini\'o la estructura de {\bf \textbackslash tableofcontents} para colocar el \'indice en una p\'agina nueva, con la palabra {\bf \'INDICE} centrada y los t\'itulos de secciones y subsecciones en tama\~no \textbackslash large.
        


        
    \subsection{Abreviaturas}
        Uso del paquete {\bf nomencl}. Los s\'imbolos est\'an alineados y tienen sangr\'ia, y se defini\'o el nombre de los grupos Acr\'onimos y S\'imbolos para diferenciar, por ejemplo, CNEA de $\epsilon$.\nomenclature[N]{$\epsilon$}{Letra griega para nomenclatura}


        
    \subsection{Im\'agenes y tablas}
    Se hace uso del paquete {\bf caption} para dar formato a los ep\'igrafes, y las referencias generan links por defecto en el pdf. Por ejemplo, ver como queda la \fig{fg:Logo_CNEA}
    \begin{figure}[ht]
        \centering
        \includegraphics[width=0.5\textwidth]{Logo_CNEA}
        \caption{Ejemplo de uso de figuras}
        \label{fg:Logo_CNEA}
    \end{figure}
    
    \par
    Por defecto, tambi\'en se carga el paquete {\bf subcaption}, de modo que se pueden insertar subfiguras como en la \fig{fg:subfiguras}. Para facilitar la uniformidad en las referencias, pueden usarse las macros {\bf \textbackslash fig, \textbackslash figs y \textbackslash tb}, tambi\'en definidas en MECOM.
    
    \begin{figure}[ht]
        \centering
        \begin{subfigure}[t]{0.45\textwidth}
            \centering
            \includegraphics[width=0.9\textwidth]{Logo_CNEA}
            \caption{Subfigura}
            \label{fg:sub1}
        \end{subfigure}
        \begin{subfigure}[t]{0.45\textwidth}
            \centering
            \includegraphics[width=0.9\textwidth]{Logo_CNEA}
            \caption{Subfigura}
            \label{fg:sub2}
        \end{subfigure} 
        \caption{Uso de subfiguras con {\bf subcaption}}
        \label{fg:subfiguras}       
    \end{figure}    
    
    
    \par
    Para usar las tablas se definieron los anchos de l\'inea con el paquete {\bf booktabs}, y pueden quedar como la \tb{tab:prueba}.
    \begin{table}[ht]
        \centering
        \begin{tabular}{c c}
            \toprule
            \bf Procesador & \bf Caracter\'istica \\
            \midrule
            \LaTeX{} & Excelente \\
            OpenOffice & Porquer\'ia\\
            \bottomrule
        \end{tabular}
        \caption{Ejemplo de uso de tablas}
        \label{tab:prueba}
    \end{table}    
    


    
    \subsection{Referencias}
    Para generar este documento se us\'o el estilo {\bf asmems4} con Bibtex, y se incluye el archivo asmems4.bst para facilitar la primera compilaci\'on. El usuario de la clase puede emplear el formato que considere m\'as adecuado.
    \par
    La clase MECOM hace uso del paquete {\bf tocbibind}, el cual aporta flexibilidad en la generaci\'on del t\'itulo e inclusi\'on en el \'indice. Por defecto, el t\'itulo de la secci\'on {\bf Bibliograf\'ia y referencias} se encuentra numerado, aunque este efecto puede cambiarse eliminando la opci\'on {\bf numbib}.
    
    
    
    \subsection{Control de cambios}
    En esta primera versi\'on, se sugiere el uso del paquete {\bf changes} para realizar revisiones (hasta que se defina una alternativa mejor). Este paquete permite, entre otras cosas, incorporar anotaciones como en \replaced[id=mas]{LibreOffice}{Microsoft Word}. Se pueden incluir autores, colores, lista de cambios, etc.
    \par
    La utilizaci\'on de este paquete en realidad no implica llevar a cabo un control de las versiones, sino remarcar las correcciones y/o extensiones que puedan incorporar los revisores. Si se usa el paquete con la opci\'on \emph{final}, la compilaci\'on produce un pdf sin estas marcas.
    
    
    
    \subsection{Formato general}
    El resto de las caracter\'isticas generales del documento, como sangr\'ia, interlineado, formato de t\'itulo de secciones, encabezados, etc., se encuentran definidas dentro de la clase, utilizando valores similares a los empleados en el template con formato .docx. La clase MECOM est\'a comentada, de forma que puede cambiarse f\'acilmente.
    
    
    
    
    \section{Conclusiones}
    Ya no hay excusas para no documentar el trabajo usando informes t\'ecnicos.\\    
    
    
    \bibliographystyle{asmems4}
    \bibliography{MiBiblioteca}
    
    
    \newpage
    \section{Prueba de m\'argenes} \blindtext[11]    


\end{document}
